% ============================================================
%  part08.tex  |  原文 PDF p.56–60（页脚 p.55–59）
%  内容：第 6 章（续）
%    6.5 量子 BV 主方程（Quantum BV Master Equation）
%    6.6 形变–障碍理论（Deformation-Obstruction Theory）
%  说明：6.4 节的标题及其 Definition 6.9–Proposition 6.13 位于 PDF p.54–55，
%        归 part07.tex；本片段自 PDF p.56 起，开头一段是 6.4 末尾的过渡文字，
%        随后进入 6.5。故此处不重复 \section 与 keypoints。
% ============================================================

% >>> p056

于是，任何规范不变的作用量（即 $f$）都自然地延拓为经典主方程的一个解（即 $S_0$）。
这就是 BV 形式的普适性。下面来看诱导出的微分
\[
  \delta = \{S_0,-\}.
\]
把 $S_0$ 写成 $S_0 = f + J$，其中
\[
  J = c^\alpha V_\alpha^i(x)\theta_i + \frac{1}{2} f^{\gamma}_{\alpha\beta}\, c^\alpha c^\beta e_\gamma .
\]
对 $J$ 稍加思索即可看出，$\{J,-\}$ 正是
$C^\bullet(\mathfrak{g},\operatorname{PV}(X)\otimes\Sym(\mathfrak{g}[2]))$ 上的
Chevalley--Eilenberg 微分。这里 $\mathfrak{g}$ 通过 Lie 导数作用于
$\operatorname{PV}(X)$，并通过伴随表示作用于 $\Sym(\mathfrak{g}[1])$。于是
\[
  \delta = d_{\mathrm{CE}} + \{f,-\},
\]
其中 $d_{\mathrm{CE}}$ 表示对称性，而
\[
  \{f,-\} = \sum_i \partial_i f \frac{\partial}{\partial \theta_i}
\]
编码了 $f$ 的临界点，即运动方程。

% [Tutor Note: 原文此段（"Thus any gauge invariant action..."）在手稿中是斜体，
% 接 PDF p.55 上 Proposition 6.13 的证明（原文仅写 "Proof: Exercise."），可视作
% 该习题的解答性注记；译文按普通正文排出，内容不作增删。另，原文 "A few moment
% thought on J" 中 "moment" 应为复数 "moments"，属笔误，按原意译出。]

% [Tutor Note: 易错提醒——原文此句前半写 C^•(g, PV(X)⊗Sym(g[2]))，紧接着却说
% g "acts on Sym(g[1]) via adjoint representation"，两处的移位数不一致；由 6.3 节
% （PDF p.53–55）deg(e_α) = -2 可知 Sym(g[2]) 才是与 {e_α} 匹配的记号，故 "g[1]"
% 疑为 "g[2]" 之笔误。译文照录原文，请核对原书 p.55。此外 d_CE 即第 4 章
% （4.1 节）的 Chevalley–Eilenberg 微分，此处它只作用在鬼场变量 c^α 上。]

\subsection{量子 BV 主方程（Quantum BV Master Equation）}

现在转向讨论量子理论。

\begin{wdefn}{6.14}
  设 $(A,\Delta)$ 是一个 BV 代数。一个 0 次元素
  \[
    S = S_0 + \hbar S_1 + \hbar^2 S_2 + \cdots \in A[[\hbar]]
  \]
  称为满足\textbf{量子（BV）主方程}（quantum (BV) master equation, QME），如果
  \[
    \hbar\Delta S + \frac{1}{2}\{S,S\} = 0.
  \]
  这里把 $\hbar$ 视为形式变量。
\end{wdefn}

在经典极限 $\hbar\to 0$ 下，QME 的首个 $\hbar$ 阶成为
\[
  \frac{1}{2}\{S_0,S_0\} = 0,
\]
这恰是经典主方程。

% [Tutor Note: 易错预警——QME 的写法与符号约定强相关。本书取
% ℏΔS + ½{S,S} = 0；另一些文献（以及作者本人的学生笔记）写成
% {S,S} + 2ℏΔS = 0，二者只差整体乘以 2；还有作者把 Δ 定义成散度算子的相反数，
% 于是符号整体翻转。引用他人结果前务必先核对其 Δ 与 {–,–} 的约定，
% 否则 6.6 节障碍类的符号会全部错位。]

\begin{wprop}{6.15}
  量子主方程等价于
  \[
    \Delta\left(\ee^{\frac{1}{\hbar}S}\right) = 0.
  \]
\end{wprop}

\begin{proof}
  这由 $\Delta$ 的“二阶”性质推出；细节留给读者。
\end{proof}

% [Tutor Note: 此处省略的论证实质——BV 算子是二阶算子，满足
% Δ(fg) = Δf·g + (-1)^{|f|}f·Δg + (-1)^{|f|}{f,g}。代入 f = g = e^{S/ℏ} 得
% Δ(e^{S/ℏ}) = e^{S/ℏ}(ℏ^{-1}ΔS + ℏ^{-2}½{S,S})，故在形式幂级数意义下
% “QME ⟺ Δ(e^{S/ℏ}) = 0”。这一步只是把定义式乘以 e^{S/ℏ} 再整理，
% 不涉及任何额外假设，但要注意 e^{S/ℏ} 不是 A[[ℏ]] 中的元素而是形式指数，
% 展开时按 ℏ 的幂次逐项比较。]

这一表达式给出了 QME 作为量子规范一致性条件的精确物理含义。在微积分中，
若 $\omega$ 是闭微分形式 $d\omega = 0$，则积分
\[
  \int_\gamma \omega
\]
% >>> p057
在闭链 $\gamma$ 上取值，并且该值在 $\gamma$ 的连续形变下保持不变。

BV 理论中有类似的现象。若 $\Delta(\ee^{S/\hbar}) = 0$，则有一种 BV 积分的概念
\[
  \int_{\mathcal{L}} \ee^{\frac{S}{\hbar}},
\]
它在 $\mathcal{L}$ 的连续形变下保持不变。这里 $\mathcal{L}$ 是与规范固定条件相关的
超 Lagrangian 子空间。因此，若 $S$ 满足 QME，则
\[
  \int_{\mathcal{L}} \ee^{\frac{S}{\hbar}}
\]
在规范固定条件的连续变化下保持不变。这正是量子规范一致性的精确表述。

% [Tutor Note: 物理图像点拨——这里的类比是：Δ 扮演外微分 d 的角色，
% e^{S/ℏ} 扮演“闭形式”，而超 Lagrangian 子空间 L（对应一个规范固定切片）
% 扮演“闭链”。规范固定条件的连续变化对应 L 的连续形变，Δ(e^{S/ℏ}) = 0
% （即 QME）保证积分不随形变而变。需要注意，BV 积分 ∫_L 只是一个形式记号
% （见第 7 章借助 (-1)-平移辛结构才把它严格化），不要把它与第 1 章尚未定义的
% 路径积分测度混为一谈。]

假设 $S$ 满足 QME。定义量子 BRST 算子
\[
  \delta^\hbar = \hbar\Delta + \{S,-\}.
\]

% [Tutor Note: 记号提醒——量子 BRST 算子记作 δ^ℏ（ℏ 在上标），经典 BRST 算子
% 记作 δ = {S_0,-}。容易与 6.6 节的 δ_S = {S,-} 混淆：δ_S 只是括号算子本身，
% δ^ℏ = ℏΔ + δ_S。另外 δ^ℏ 一般不再是分次导子（Δ 破坏 Leibniz 法则），
% 但它满足 (δ^ℏ)^2 = 0，仍可定义上同调，此点见 Proposition 6.16。]

\begin{wprop}{6.16}
  $\delta^\hbar$ 满足 $(\delta^\hbar)^2 = 0$。
\end{wprop}

\begin{proof}
  记 $\delta_S = \{S,-\}$。利用
  \[
    \Delta\{S,-\} = -\{\Delta S,-\} - \{S,\Delta(-)\},
  \]
  可得
  \[
    \Delta\delta_S + \delta_S\Delta = -\{\Delta S,-\}.
  \]
  于是
  \[
    \begin{aligned}
      (\delta^\hbar)^2 &= (\hbar\Delta)^2 + \hbar(\Delta\delta_S + \delta_S\Delta) + \delta_S^2\\
      &= -\{\hbar\Delta S,-\} - \frac{1}{2}\{\{S,S\},-\}\\
      &= -\left\{\hbar\Delta S + \frac{1}{2}\{S,S\},-\right\} = 0.
    \end{aligned}
  \]
\end{proof}

% [Tutor Note: 符号约定警示——第二行到第三行之间用了 Jacobi 恒等式把 δ_S²
% 化为 ±½{{S,S},-}；这里出现的负号与本书“Δ{S,-} = -{ΔS,-} - {S,Δ(-)}”
% 的约定绑定。若改用相反的 BV 括号符号（或把 Δ 换成 -Δ），整段的符号会
% 一起翻转，但结论 (δ^ℏ)² = 0 不变。核对时请以定义式（QME 与 Δ 的二阶性质）
% 为准，不要只凭记忆套用；这是 BV 形式中最常见的符号陷阱。]

显然 $\lim_{\hbar\to0}\delta^\hbar = \delta = \{S_0,-\}$ 给出经典 BRST 算子。
因此 $\delta^\hbar$ 可视为 $\delta$ 的量子化。微分 $\delta$ 与 $\delta^\hbar$
分别刻画经典与量子的规范变换。

\begin{wdefn}{6.17}
  设 $(A,\Delta)$ 为 BV 代数，并假设 $S = S_0 + \hbar S_1 + \cdots$ 解量子主方程
  （从而 $S_0$ 解经典主方程）。
  \begin{enumerate}[label=\textcircled{\small\arabic*},leftmargin=2.4em,itemsep=0.2em]
    \item 上同调
          \[
            \Obs^c = H^\bullet(A,\{S_0,-\})
          \]
          称为\textbf{经典可观测量}空间。
    \item 上同调
          \[
            \Obs^q = H^\bullet(A[[\hbar]],\hbar\Delta+\{S,-\})
          \]
          称为\textbf{量子可观测量}空间。
  \end{enumerate}
  它们表示模去规范等价后的经典与量子规范不变对象。
\end{wdefn}

% [Tutor Note: 记号提醒——H^• 中的 • 是“分级”占位符，意为对所有次数的上同调
% 取直和；Obs^c 的上标 c 指 classical，Obs^q 的 q 指 quantum，不是指数。
% 物理上，H^0 对应规范不变的可观测量，H^1 对应无穷小形变/障碍（见 6.6 节），
% 更高次上同调则对应高阶规范结构。]

% >>> p058

\begin{wexam}{6.18}（奇点理论 / Landau--Ginzburg B-模型）
  设 $f:\C^n\to\C$ 是一个在原点有孤立临界点的多项式。考虑由多项式全纯多向量场
  构成的 BV 代数 $(A,\Delta)$：
  \[
    A = \C[z^i,\theta_i],\qquad \theta_i\theta_j = -\theta_j\theta_i,\qquad
    \Delta = \sum_i \frac{\partial}{\partial z^i}\frac{\partial}{\partial \theta_i}.
  \]
  则
  \begin{itemize}[leftmargin=1.6em,itemsep=0.25em]
    \item $f = f(z^i)$ 满足 QME，因为 $f$ 不依赖 $\theta_i$：
      \[
        \Delta\left(\ee^{\frac{1}{\hbar}f(z)}\right) = 0.
      \]
    \item 经典可观测量：
      \[
        \Obs^c = H^\bullet(A,\{f,-\}) = \operatorname{Jac}(f)
        = \C[z^i]/(\partial_i f).
      \]
      这里 $\operatorname{Jac}(f)$ 是 Jacobi 环，$(\partial_i f)$ 是由
      $\partial_i f$ 生成的理想。
    \item 量子可观测量：
      \[
        \Obs^q = H^\bullet(A[[\hbar]],\hbar\Delta+\{f,-\})
        \simeq \Omega^n[[\hbar]]/(\hbar\,\dd + \dd f\wedge)\,\Omega^{n-1}[[\hbar]].
      \]
      这称为（形式的）Brieskorn 格（Brieskorn lattice），在奇点理论中起重要作用。
    \item BV 积分模拟振荡积分
      \[
        \int_{\mathcal{L}} \ee^{\frac{1}{\hbar}f}\mathcal{O},
      \]
      其中 $\mathcal{L}$ 是一个 Lefschetz 快速下降链（Lefschetz thimble）。
  \end{itemize}
  Brieskorn 格与振荡积分都与 Hodge 结构深刻相关。事实上，$\hbar$-进滤过恰是
  Hodge 滤过。因此量子场论可以看作 $\infty$ 维 Hodge 理论的一个版本。
\end{wexam}

% [Tutor Note: 术语与记号依据——（i）本例中出现的花体 \mathcal{O}（BV 积分里的
% 被积可观测量，以及 6.6 节中的障碍类 \mathcal{O}_k）与渐近分析的大 O 记号
% O(ℏ^k) 形状相近但含义完全不同，切勿混用；前者应排作 \mathcal{O}。
% （ii）Jac(f) = C[z^i]/(∂_i f) 称为 f 的 Jacobi 环（局部代数），其维数等于
% f 的孤立临界点重数（Milnor 数）。（iii）“Brieskorn lattice” 中 lattice 指
% D-模意义下的格（与 Hodge 滤过相伴的 Gauss–Manin 联络的格子），故译
% “Brieskorn 格”，不译“晶格”。Lefschetz thimble 是 Picard–Lefschetz 理论中
% 从临界点出发的快速下降流形，译“Lefschetz 快速下降链”。原文“the ℏ-adic
% filtration is precisely the Hodge filtration”一句是本书的核心洞见之一。]

\subsection{形变–障碍理论（Deformation-Obstruction Theory）}

现在我们来回答如下自然的问题：从经典主方程的一个解 $S_0$ 出发，是否总能把它
量子化为量子主方程的解 $S = S_0 + O(\hbar)$？

一般而言，答案是\textbf{否}。可观测量复形恰好控制这一形变理论，并捕捉其中的障碍。

设 $(A,\Delta)$ 为 BV 代数。假设 $S_0$ 解经典主方程，且 $\delta = \{S_0,-\}$。
我们要找 $S_1,S_2,\cdots\in A$，使得
\[
  S = S_0 + S_1\hbar + S_2\hbar^2 + \cdots
\]
解量子主方程。我们按 $\hbar$ 的阶逐阶分析。

% [Tutor Note: 这一段是本节的总纲：把“量子化”看成一个形变问题——在经典解
% S_0 的基础上逐阶添加 ℏ 修正。H^1(A,δ) 度量“下一步是否走得通”，因此它被称为
% 障碍空间；H^0(A,δ) 度量“走通之后解有多少自由度”，因此被称为切空间/形变空间。
% 这一图景就是标准的形变–障碍理论（Kodaira–Spencer/Kuranishi 型），
% 只不过这里的“复形”是 BV 可观测量复形。]

% >>> p059

\paragraph{$\hbar^1$ 阶}

QME 的 $\hbar^1$ 项为
\[
  \hbar\Delta S_0 + \hbar\{S_0,S_1\} = 0,
\]
即我们要找解 $S_1$ 满足
\[
  \delta S_1 = -\Delta S_0.
\]
由于 $\delta^2 = 0$，$S_1$ 可解的一个必要条件是
\[
  \delta(\Delta S_0) = 0.
\]
这确实成立，因为由经典主方程有
\[
  \{S_0,\Delta S_0\} = -\frac{1}{2}\Delta\{S_0,S_0\} = 0.
\]
令 $\mathcal{O}_1 = \Delta S_0$。由于
\[
  \delta\mathcal{O}_1 = 0,
\]
$\mathcal{O}_1$ 代表一个上同调类 $[\mathcal{O}_1]\in H^1(A,\delta)$。上述方程表明，
$S_1$ 可解当且仅当 $\mathcal{O}_1$ 是 $\delta$-恰当的。

% [Tutor Note: 易错预警——关系式 {S_0, ΔS_0} = -½Δ{S_0,S_0} 来自
% “Δ 把 BV 括号变成括号的导子”这一恒等式；使用它时符号取决于 Δ 与 {–,–}
% 的约定（参见 Proposition 6.16 证明后的批注）。此处右侧为零是因为
% {S_0,S_0} = 0（经典主方程），而 Δ0 = 0。切记 S_0 是 0 次（偶）元素，
% 因此不需要额外的 Koszul 符号；若把 S_0 误当奇元素，符号会出错。]

\begin{wprop}{6.19}
  $S_1$ 可解当且仅当 $H^1(A,\delta)$ 中 $[\mathcal{O}_1] = 0$。
\end{wprop}

现在假设 $[\mathcal{O}_1] = 0$。设 $S_1$ 与 $\widetilde{S}_1$ 是 $\hbar^1$ 阶的
两个解。则
\[
  \delta S_1 = \delta\widetilde{S}_1 = -\mathcal{O}_1,
\]
因此 $\widetilde{S}^1 - S^1$ 是 $\delta$-闭的。另一方面，对任意 $J\in A_0$，
解 $S_1 + \delta J$ 与 $S_1$ 在适当意义下是规范等价的。于是，$S_1$ 的解模去
规范等价后为
\[
  \begin{cases}
    \varnothing, & \text{if } [\mathcal{O}_1]\neq 0,\\
    H^0(A,\delta), & \text{if } [\mathcal{O}_1] = 0.
  \end{cases}
\]
$\delta$-上同调类 $[\mathcal{O}_1]$ 称为 $\hbar^1$ 阶的\textbf{障碍类}
（obstruction class）。

% [Tutor Note: 原文疑误两处，均照录原文、请核对原书 p.58–59：
% （i）“represents a cohomology clas” 中 “clas” 应为 “class”；
% （ii）“thus S̃¹ - S¹ is δ-closed” 的上标 1 疑为下标之误，按上下文应为
% S̃_1 - S_1。译文保留原文的上下标位置，并在正文中照录，以便读者对照原书。]

\paragraph{$\hbar^{>1}$ 阶}

假设已找到
\[
  S_{<k} := S_0 + \hbar S_1 + \cdots + \hbar^{k-1}S_{k-1},
\]
它把 QME 解到 $\hbar^{k-1}$ 阶：
\[
  \hbar\Delta S_{<k} + \frac{1}{2}\{S_{<k},S_{<k}\} = O(\hbar^k).
\]
下面考虑求解 $S_k$ 的问题。上述方程可等价地写成
\[
  \Delta\ee^{\frac{1}{\hbar}S_{<k}} = O(\hbar^{k-2})\ee^{\frac{1}{\hbar}S_{<k}}.
\]
记
\[
  \Delta\ee^{\frac{S_{<k}}{\hbar}}
  = \left(\hbar^{k-2}\mathcal{O}_k + O(\hbar^{k-1})\right)\ee^{\frac{S_{<k}}{\hbar}},
\]
其中 $\mathcal{O}_k$ 是 $O(\hbar^{k-2})$ 中的首项。等价地，我们有
\[
  \hbar\Delta S_{<k} + \frac{1}{2}\{S_{<k},S_{<k}\}
  = \hbar^k\mathcal{O}_k + O(\hbar^{k+1}).
\]

% [Tutor Note: 省略论证点拨——从 Δe^{S_{<k}/ℏ} 过渡到 ℏΔS_{<k} + ½{S_{<k},S_{<k}}
% 仍用 Δ 的二阶性质：两者相差一个因子 e^{S_{<k}/ℏ} 与 ℏ^{-2}（见 Proposition 6.15
% 后的批注），故 Δe^{S/ℏ} 的 ℏ^{k-2} 阶首项恰好对应 QME 表达式的 ℏ^k 阶首项。
% 注意此处的展开默认 k ≥ 2；k = 1 的情形已在上一段单独处理，不能直接套用
% ℏ^{k-2} 的幂次。]

% >>> p060

\begin{wclaim}{}
  $\mathcal{O}_k$ 是 $\delta$-闭的。
\end{wclaim}

\begin{proof}
  对 $\mathcal{O}_k$ 的定义式作用 $\Delta$，得到
  \[
    \begin{aligned}
      \Delta\left(\left(\hbar^{k-2}\mathcal{O}_k + O(\hbar^{k-1})\right)
        \ee^{\frac{1}{\hbar}S_{<k}}\right)
        &= \Delta^2 \ee^{\frac{1}{\hbar}S_{<k}} = 0\\
      \Longrightarrow\quad
      \hbar^{k-3}\left(\delta\mathcal{O}_k + O(\hbar)\right)
        \ee^{\frac{1}{\hbar}S_{<k}} &= 0\\
      \Longrightarrow\quad \delta\mathcal{O}_k &= 0.
    \end{aligned}
  \]
\end{proof}

% [Tutor Note: 幂次核对——对 ℏ^{k-2}𝒪_k·e^{S_{<k}/ℏ} 作用 Δ 时，Δ 作用在
% 指数因子上会多出 ℏ^{-1}，故首项为 ℏ^{k-3}δ𝒪_k；这正是第二行出现 ℏ^{k-3}
% 的原因（k ≥ 3 时该项主导）。原文 Claim 不编号，故本译本用无编号的“断言”
% 环境；若按编号环境排版会与原文编号体系不符。]

求解 $S_k$ 的问题在于把 QME 解到 $\hbar^k$ 阶：
\[
  \hbar\Delta(S_{<k}+\hbar^k S_k)
  + \frac{1}{2}\{S_{<k}+\hbar^k S_k,\ S_{<k}+\hbar^k S_k\}
  = O(\hbar^{k+1}).
\]
这等价于
\[
  \hbar^k\{S_0,S_k\} + \hbar^k\mathcal{O}_k = 0,
\]
即
\[
  \delta S_k = -\mathcal{O}_k.
\]
我们处于与先前相同的情形。$S_k$ 是否可解，等同于问 $\mathcal{O}_k$ 是否
$\delta$-恰当。于是，$S_k$ 的解模去规范等价后为
\[
  \begin{cases}
    \varnothing, & \text{if } [\mathcal{O}_k]\neq 0,\\
    H^0(A,\delta), & \text{if } [\mathcal{O}_k] = 0.
  \end{cases}
\]
这里 $[\mathcal{O}_k]\in H^1(A,\delta)$ 是把 QME 解到 $\hbar^k$ 阶的
\textbf{障碍类}（规范反常）。空间 $H^1(A,\delta)$ 是\textbf{障碍空间}，而
$H^0(A,\delta)$ 是\textbf{切空间}或\textbf{形变空间}。上述分析给出下列关于
构造量子化的“抽象废话”（abstract nonsense）的有力定理。

\begin{wthm}{6.20}
  若 $H^1(A,\delta) = 0$，则 $S_0$ 可以量子化为 QME 的一个解
  $S = S_0 + \hbar S_1 + \cdots$。
\end{wthm}

% [Tutor Note: 术语与物理含义——“abstract nonsense” 是范畴论/同调代数的自嘲式
% 行话，指“用抽象结构（复形、上同调、障碍类）而非具体计算就能得到的普遍结论”，
% 译“抽象废话”并保留英文，以免读者误以为作者在贬低该定理。物理上，[𝒪_k] 即
% 通常所说的“规范反常”（gauge anomaly）：若 H^1(A,δ) ≠ 0，则经典作用量 S_0
% 无法逐阶量子化，微扰量子化在 ℏ^k 阶被障碍类卡住。定理 6.20 的条件 H^1 = 0
% 是充分条件而非必要条件：即使 H^1 ≠ 0，只要所有障碍类 [𝒪_k] 恰好为零，
% 量子化仍可进行。此定理只保证形式幂级数解的存在性，不保证其收敛性
% （后者属于 Borel 可和性等分析问题）。]
